\section*{Phase 4: Modeling}

\subsection*{Task 4.1 Selecting Modeling Techniques}

\subsubsection*{Defined Modeling Technique(s)}

\subsubsection*{Comparative Analysis of Influenza Smoothing Techniques}

We implement wavelet transformations for the denoising of some of our features. Through the application of smoothing techniques: Exponentially Weighted Moving Average (EWA), Moving Average (MA), and Discrete Wavelet Transformation (DWT), distinct differences can be seen in the performance of capturing the trends for influenza cases. Specifically, DWT demonstrates the greatest effectiveness in noise reduction while still capturing important trend inflection points, which speaks to his effectiveness in preserving meaningful patterns without sacrificing clarity. In contrast, the 5-day MA results in a smoother curve; however, there is considerable lag, causing it to fail to capture swift outbreak surges in a timely manner. The EWA balances these two extremes since it adapts quicker than MA by emphasizing recent data, although it does display slightly more variability than DWT. Ultimately, while DWT is most effective for retaining signal dynamics, EWA is best for responsive monitoring, and MA is best for long-term analysis. Below are the Discrete Wavelet Transformation (DWT):

\begin{figure}[htbp]
\centering
\includegraphics[width=\textwidth]{10_wavelet_denoising.png}
\caption{Discrete Wavelet Transformation (DWT) denoising.}
\end{figure}

Moving Average (MA):

\begin{figure}[htbp]
\centering
\includegraphics[width=\textwidth]{11_ma_comparison.png}
\caption{Moving average (MA) comparison.}
\end{figure}

and Exponentially Weighted Moving Average (EWA):

\begin{figure}[htbp]
\centering
\includegraphics[width=\textwidth]{12_ewa_comparison.png}
\caption{Exponentially weighted moving average (EWA) comparison.}
\end{figure}

To enhance both the accuracy of the predictions and their interpretability, we selected three modeling approaches:

\begin{enumerate}
  \item \textbf{Linear Regression (LR)}
  \begin{itemize}
    \item Considered a baseline model because of its straightforward nature and ease of understanding.
    \item Helps evaluate the presence of linear relationships between the chosen features and influenza case volumes.
    \item In the case that interactions between the inputs and outputs are continual and linear over time, LR is expected to result in the best performance.
    \item Often struggles to accommodate cyclical seasonal effects or sharp trend breaks, which are frequently present in epidemic data.
  \end{itemize}
  \item \textbf{Decision Tree Regressor (DTR)}
  \begin{itemize}
    \item Best suited for capturing non-linear patterns dictating influenza outbreak cycles.
    \item Simple to interpret and presents feature importance in an easy to understand manner.
    \item Robust to outliers and does not require feature scaling, making this model particularly useful for noisy epidemiological data.
    \item May overfit the data if not constrained (through methods such as pruning or limiting depth); this is especially true when trained on small or high-variance datasets.
  \end{itemize}
  \item \textbf{Neural Network Regressor (Multilayer Perceptron, MLP)}
  \begin{itemize}
    \item Learns complex, non-linear relationships between features.
    \item Used to model the percentage change in trend volume in influenza cases.
    \item Applies dropout and batch normalization to mitigate overfitting and improve generalization.
    \item Is slower to train, requiring more hyperparameter optimization, and may need more data to outperform simpler models.
  \end{itemize}
\end{enumerate}

All models described were trained using two feature sets.: \textbf{\texttt{All\_Features}} ('\texttt{ISO\_YEAR}', '\texttt{ISO\_WEEK}', '\texttt{Month}', '\texttt{INF\_ALL\_LAG\_1}', '\texttt{INF\_ALL\_LAG\_2}', '\texttt{INF\_ALL\_LAG\_4}', '\texttt{INF\_A}', '\texttt{INF\_B}', '\texttt{ILI\_ACTIVITY}', '\texttt{INF\_ALL}', '\texttt{SPEC\_PROCESSED\_NB}', '\texttt{Total\_Influenza\_Cases}') and \textbf{\texttt{Season\_Features}} ('\texttt{Season\_numeric}', '\texttt{INF\_ALL}', '\texttt{SPEC\_PROCESSED\_NB}', '\texttt{Total\_Influenza\_Cases}'), allowing comparative performance assessment.

\subsubsection*{Defined Modeling Assumptions}

\begin{itemize}
  \item \textbf{Feature Relevance and Sufficiency}: We heuristically (which could be a bias) selected the features such as past flu cases, temperature, and population density to build a model to predict future outbreaks of flu.
  \item \textbf{Data Quality}: We hypothesize that the historical data (such as WHO data, etc.) is valid and will not have any discrepancies or gaps.
  \item \textbf{No Structural Breaks}: We do not make assumptions about disruptive changes, for example, new data-busting strains like COVID-19.
  \item \textbf{Temporal Structure}: Given that the influenza dataset has a high value temporal dependency, lag features (\texttt{INF\_ALL\_LAG\_1}, \texttt{\_2}, \texttt{\_4}) were added to reflect the trends of the past.
  \item \textbf{Stationarity Assumption}: Stabilized variance and learning by using wavelet denoising and percentage change transformations.
  \item \textbf{Independence of Observations}: While dependencies in time order are defined, we consider feature vectors at any given point in time to be independent after certain transformations for training the model to be independent.
  \item \textbf{Feature Normalization}: Utilized a StandardScaler for setting a uniform measurement for all features to be used in neural networks and tree-based models.
  \item \textbf{Model Assumptions}:
  \begin{itemize}
    \item \textbf{Linear Regression}
    \begin{itemize}
      \item Linearity: The relation of features to the target variable is assumed to be linear.
      \item No Multicollinearity: Features are supposed to be uncorrelated or have low correlations with each other.
      \item Homoscedasticity: The variance of the errors (residuals) must be constant across all levels of the input variables.
      \item Independence of Errors: Residuals must not correlate with each other (especially critical in time series data).
      \item Normality of Errors: Assumed normally distributed residuals (mostly for inference, like CIs).
      \item Correct Model Specification: Assumes all relevant variables are present and exclusively irrelevant variables are excluded.
    \end{itemize}
    \item \textbf{Decision Tree Regressor}
    \begin{itemize}
      \item No Assumption of Linearity: It deals with non-linear relationships as well.
      \item No Need for Feature Scaling: Standardization or normalization is not required.
      \item No Assumption on Distribution of Errors: It does not make assumptions of homoscedasticity or normal distribution of residuals.
      \item Handles Multicollinearity Better: While multicollinearity does influence interpretability, trees tend to select one dominant feature and ignore the rest.
      \item No Assumption of Independence: Compared to regression or neural nets, trees do not require iid samples for the purpose of independence.
    \end{itemize}
    \item \textbf{Neural Network Regressor}
    \begin{itemize}
      \item Large, Representative Data: Neural networks expect you possess enough data that patterns are present without overfitting.
      \item Features Should Be Scaled: Inputs being normalized or standardized yields better results during the training phase.
      \item Independent and Identically Distributed: Training examples come from the same distribution and independent of each other, which makes the model learn better from available examples.
      \item Low Noise-to-Signal Ratio: Excessive noise or irrelevant data can lead to obscuring important information and differentiating signals.
      \item Stationarity for Temporal Patterns: The model accepts working with time series data like flu cases over time on the assumption that certain patterns tend to remain in the long-run exception is when using RNNs or transformers which model time dynamically.
    \end{itemize}
  \end{itemize}
\end{itemize}

\subsection*{Task 4.3 Building Model(s)}

\subsubsection*{Parameter Definitions}

With this project, model building included three types of models: \textbf{Linear Regression, Decision Tree Regressor, and a Neural Network (MLP)}. Each model required its own unique configuration; some configurations were preset, while others were tailored:

\begin{itemize}
  \item \textbf{Linear Regression (Scikit-learn):}
  \begin{itemize}
    \item Parameters used: All defaults (fit\_intercept=True, normalize='deprecated', n\_jobs=None, etc.)
    \item Rationale: Linear regression was kept simple to serve as a baseline model.
  \end{itemize}
  \item \textbf{Decision Tree Regressor (Scikit-learn):}
  \begin{itemize}
    \item Parameters used: All defaults (criterion='squared\_error', max\_depth=None, min\_samples\_split=2, etc.)
    \item Rationale: Default settings were retained to evaluate initial tree performance. Tuning was deferred for initial model comparison.
  \end{itemize}
  \item \textbf{Neural Network (Keras MLP):}
  \begin{itemize}
    \item Architecture:
    \begin{itemize}
      \item Flatten $\rightarrow$ BatchNorm $\rightarrow$ Dense(64, elu) $\rightarrow$ BatchNorm $\rightarrow$ Dropout(0.2)
      \item Dense(32, elu) $\rightarrow$ BatchNorm $\rightarrow$ Dropout(0.2)
      \item Dense(16, elu) $\rightarrow$ BatchNorm $\rightarrow$ Dropout(0.2) $\rightarrow$ Output(Dense(1))
    \end{itemize}
    \item Optimizer: RMSprop with learning\_rate=0.001
    \item Loss function: MAE
    \item Callbacks: EarlyStopping (patience=10), ModelCheckpoint (best weights saved)
    \item Epochs: 120
    \item Input shape: Reshaped to (samples, 1, features) for Keras compatibility
    \item Rationale: The architecture was chosen to provide sufficient depth while ensuring overfitting is mitigated, and is with ELU activation functions for improved gradient flow.
    \item Data Preprocessing:
    \begin{itemize}
      \item StandardScaler was applied to selected numerical columns to normalize scales across models.
      \item Lag features (1, 2, 4 weeks) were created to support time-dependent forecasting.
      \item Outliers were removed using IQR to improve generalization.
      \item Wavelet Denoising was applied to smooth time-series trends and compared with MA/EMA smoothing.
      \item Categorical Feature (Season) was encoded numerically.
    \end{itemize}
  \end{itemize}
\end{itemize}

All parameter choices and preprocessing steps have been documented and justified based on initial model performance and literature-supported practices.

\subsubsection*{Model Description}

\begin{itemize}
  \item \textbf{Linear Regression}:
  \begin{itemize}
    \item Why: It serves as a benchmark for estimating trends.
    \item Interpretation: Each coefficient quantifies the impact of an input feature on the number of predicted influenza cases.
    \item Variables: Used time features (year, week, month), influenza strain counts, lagged values, and activity index.
    \item Difficulties: Issues with multicollinearity as a strong negative effect; while there are mitigations, lag features do little to solve autocorrelation concerns.
  \end{itemize}
  \item \textbf{Decision Tree}:
  \begin{itemize}
    \item Why: Non-parametric model well-suited for capturing non-linear relations without the need for feature standardization.
    \item Interpretation: Rules of decision are visual with post evaluation on change in performance satisfactory in target metric driven approach consigned.
    \item Variables: Like above but applied on percentage change data.
    \item Difficulties: Pruning and setting maximum depth limits helps reduce overfitting but are needed.
  \end{itemize}
  \item \textbf{Neural Network}:
  \begin{itemize}
    \item Why: Captures sophisticated non-linear relationships while overfitting can be controlled for with dropout regularization and batch normalization.
    \item Interpretation: Model viewed as a black-box interpretable through given performance metrics (MAE, MSE).
    \item Variables: Percentage change inputs excluding the target variable across two different sets of features were standardized and incorporated.
    \item Difficulties: Greater training durations along with an increased sensitivity to imbalance in the data, hyperparameter setting, learning rate, and rate of adjustment made for dependent variables, and the chosen method of shifting.
  \end{itemize}
\end{itemize}

\subsubsection*{Feature Sets Evaluated:}

\begin{itemize}
  \item \textbf{\texttt{All\_Features}} ('\texttt{ISO\_YEAR}', '\texttt{ISO\_WEEK}', '\texttt{Month}', '\texttt{INF\_ALL\_LAG\_1}', '\texttt{INF\_ALL\_LAG\_2}', '\texttt{INF\_ALL\_LAG\_4}', '\texttt{INF\_A}', '\texttt{INF\_B}', '\texttt{ILI\_ACTIVITY}', '\texttt{INF\_ALL}', '\texttt{SPEC\_PROCESSED\_NB}', '\texttt{Total\_Influenza\_Cases}'): Comprehensive including lags and strain types.
  \item \textbf{\texttt{Season\_Features}} ('\texttt{Season\_numeric}', '\texttt{INF\_ALL}', '\texttt{SPEC\_PROCESSED\_NB}', '\texttt{Total\_Influenza\_Cases}'): Seasonal signal encoding and core indicators.
  \item Both with and without target variable included, for feature selection/testing impact.
\end{itemize}

\subsubsection*{Data Models}

\begin{itemize}
  \item Models were built using features derived from:
  \begin{itemize}
    \item Seasonal categorization (\texttt{Season\_numeric}), lag variables, and processed case counts.
    \item Two datasets: Full set (\texttt{All\_Features}) and a reduced seasonal set (\texttt{Season\_Features}).
  \end{itemize}
  \item Data splits were made chronologically into:
  \begin{itemize}
    \item Training set (1997--2019), validation set (2020--2022), and test set (2023--2024).
  \end{itemize}
  \item Data transformations included:
  \begin{itemize}
    \item Lag creation, wavelet denoising, moving averages, outlier removal (IQR method), and scaling.
  \end{itemize}
\end{itemize}

\subsection*{Task 4.4 Assessing Model(s)}

\subsubsection*{Model Assessment}

Models were evaluated on the validation set using \textbf{MAE} and \textbf{MSE}:
\begin{itemize}
  \item \textbf{MAE} was used for interpretability of absolute prediction error.
  \item \textbf{MSE} emphasized penalizing larger errors.
\end{itemize}

\begin{longtable}{|>{\raggedright\arraybackslash}p{3cm}|>{\raggedright\arraybackslash}p{4cm}|>{\raggedright\arraybackslash}p{3cm}|>{\raggedright\arraybackslash}p{3cm}|}
\hline
Feature Set & Model & MAE & MSE \\
\hline
\endhead
\hline
All Features & Simple Linear Regression & 45.763969 & 3,981.419 \\
\hline
All Features & Simple Decision Tree & 0.774098 & 2.062631 \\
\hline
All Features & Neural Network Regressor & 0.390836 & 1.662638 \\
\hline
Seasonal Features & Simple Linear Regression & 2,061.554912 & 9,204,389 \\
\hline
Seasonal Features & Simple Decision Tree & 11.167341 & 14,898.67 \\
\hline
Seasonal Features & Neural Network Regressor & 0.492063 & 1.710484 \\
\hline
\end{longtable}

\begin{itemize}
  \item Performance was visualized in:
  \begin{itemize}
    \item \emph{MAE/MSE bar charts} comparing models across feature sets.
  \end{itemize}
\end{itemize}

\begin{figure}[htbp]
\centering
\includegraphics[width=\textwidth]{21_mae_comparison.png}
\caption{MAE comparison across models and feature sets.}
\end{figure}

\begin{figure}[htbp]
\centering
\includegraphics[width=\textwidth]{22_mse_comparison.png}
\caption{MSE comparison across models and feature sets.}
\end{figure}

\begin{itemize}
  \item \emph{Line plots} of true vs.\ predicted values for each model.
\end{itemize}

\begin{figure}[htbp]
\centering
\includegraphics[width=\textwidth]{14_lr_predictions_All_Features.png}
\caption{Linear Regression predictions (All Features).}
\end{figure}

\begin{figure}[htbp]
\centering
\includegraphics[width=\textwidth]{15_dt_predictions_All_Features.png}
\caption{Decision Tree predictions (All Features).}
\end{figure}

\begin{figure}[htbp]
\centering
\includegraphics[width=\textwidth]{16_nn_predictions_All_Features.png}
\caption{Neural Network predictions (All Features).}
\end{figure}

\begin{figure}[htbp]
\centering
\includegraphics[width=\textwidth]{18_lr_predictions_Season_Features.png}
\caption{Linear Regression predictions (Season Features).}
\end{figure}

\begin{figure}[htbp]
\centering
\includegraphics[width=\textwidth]{19_dt_predictions_Season_Features.png}
\caption{Decision Tree predictions (Season Features).}
\end{figure}

\begin{figure}[htbp]
\centering
\includegraphics[width=\textwidth]{20_nn_predictions_Season_Features.png}
\caption{Neural Network predictions (Season Features).}
\end{figure}

\textbf{Insights:}
\begin{itemize}
  \item Neural networks performed well on percentage change prediction but had higher complexity.
  \item Decision trees underperformed with not reasonable MAE.
  \item Linear regression performed exceptionally well compared with other models.
\end{itemize}
